%! Multiplying \& Dividing Radicals; Rationalizing — Unit 5, Lesson 5.3 — Homework
\documentclass[10pt]{article}
\usepackage{algebra2-article}
\usepackage{algebra2-boxes}
\newcommand{\TallMath}[1]{$\displaystyle #1\rule[-1.4em]{0pt}{3.2em}$}

\begin{document}

\pageheader{Unit 5, Lesson 5.3}{Homework}
\namedateperiod

\vspace{0.10in}

\begin{notesbox}{Practice}
Assume all variables represent \textbf{positive} real numbers. Leave no radical in any denominator.
No calculator except where a problem asks for a decimal.

\begin{enumerate}[leftmargin=1.9em, itemsep=11pt]
  \item \textbf{Multiply and simplify.} Remember: the product is where a perfect power often
        appears out of nowhere, so finish every one.
    \begin{center}
    \renewcommand{\arraystretch}{1.9}
    \begin{tabularx}{\linewidth}{@{}X X X X@{}}
    (a) \; $\sqrt5\cdot\sqrt{20}=$ \blank{1.7cm} & (b) \; $\sqrt6\cdot\sqrt{14}=$ \blank{1.7cm} &
    (c) \; $4\sqrt3\cdot2\sqrt6=$ \blank{1.7cm} & (d) \; $\left(3\sqrt7\right)^{2}=$ \blank{1.7cm} \\
    (e) \; $\sqrt[3]{4}\cdot\sqrt[3]{10}=$ \blank{1.7cm} &
    (f) \; $\sqrt{6x}\cdot\sqrt{3x^{3}}=$ \blank{1.9cm} &
    (g) \; $\sqrt{10}\cdot\sqrt{10}=$ \blank{1.7cm} &
    (h) \; $2\sqrt{12}\cdot\sqrt{27}=$ \blank{1.7cm} \\
    \end{tabularx}
    \end{center}

  \item \textbf{Expand and simplify.} Distribute or FOIL, then collect like radicals.
    \begin{center}
    \renewcommand{\arraystretch}{2.1}
    \begin{tabularx}{\linewidth}{@{}X X@{}}
    (a) \; $\sqrt2\left(\sqrt8+\sqrt{18}\right)=$ \blank{2.6cm} &
    (b) \; $\sqrt5\left(3-\sqrt{15}\right)=$ \blank{2.6cm} \\
    (c) \; $\left(1+\sqrt3\right)\left(4-\sqrt3\right)=$ \blank{2.6cm} &
    (d) \; $\left(\sqrt6-\sqrt2\right)^{2}=$ \blank{2.6cm} \\
    (e) \; $\left(5+\sqrt7\right)\left(5-\sqrt7\right)=$ \blank{2.6cm} &
    (f) \; $\left(2\sqrt3+\sqrt5\right)\left(2\sqrt3-\sqrt5\right)=$ \blank{2.6cm} \\
    \end{tabularx}
    \end{center}
    Two of the six are \textbf{conjugate pairs}, and you could tell before multiplying. Which two?
    \blank{2.4cm} \; A \emph{third} problem also came out with no radical, for a completely
    different reason. Which one, and what was the reason? \writeline

  \item \textbf{Divide, and rationalize where you must.}
    \begin{center}
    \renewcommand{\arraystretch}{2.3}
    \begin{tabularx}{\linewidth}{@{}X X@{}}
    (a) \; $\dfrac{\sqrt{72}}{\sqrt{2}}=$ \blank{2.2cm} &
    (b) \; $\dfrac{\sqrt{45x^{3}}}{\sqrt{5x}}=$ \blank{2.2cm} \\
    (c) \; $\dfrac{7}{\sqrt{2}}=$ \blank{2.2cm} &
    (d) \; $\dfrac{\sqrt{3}}{\sqrt{8}}=$ \blank{2.2cm} \\
    (e) \; $\sqrt{\dfrac{2}{7}}=$ \blank{2.2cm} &
    (f) \; $\dfrac{10}{4+\sqrt{6}}=$ \blank{2.6cm} \\
    (g) \; $\dfrac{\sqrt{2}}{\sqrt{5}+\sqrt{3}}=$ \blank{2.6cm} &
    (h) \; $\dfrac{6}{\sqrt{3}-3}=$ \blank{2.6cm} \\
    \end{tabularx}
    \end{center}
    \emph{Watch the sign in (h).}

  \item \textbf{Two questions about why this works.}
    \begin{enumerate}[label=(\alph*), leftmargin=1.8em, itemsep=6pt]
      \item A classmate claims $\left(\sqrt3+\sqrt5\right)^{2}=3+5=8$. Use decimals to show this is
            wrong, then expand it correctly and name the term they threw away.
            \writelines{3}
      \item To rationalize $\dfrac{1}{2+\sqrt3}$, why does multiplying top and bottom by
            $2+\sqrt3$ fail while $2-\sqrt3$ succeeds? And why is either one \emph{allowed} in the
            first place --- what keeps the value of the expression from changing?
            \writelines{3}
    \end{enumerate}

  \item \textbf{Two roads, one answer (connecting to Lesson 5.1).} Rationalize
        $\dfrac{1}{\sqrt[3]{5}}$ two ways.
    \begin{enumerate}[label=(\alph*), leftmargin=1.8em, itemsep=5pt]
      \item \emph{With radicals.} You need \blank{0.8cm} copies of $5$ under the cube root, so
            multiply top and bottom by \blank{1.6cm}: \; $\dfrac{1}{\sqrt[3]5}=\blank{2.2cm}$
      \item \emph{With exponents.} $\dfrac{1}{\sqrt[3]5}=5^{\blank{1.0cm}}
            =\dfrac{5^{\blank{1.0cm}}}{5}=\blank{2.2cm}$
      \item Why must the two roads agree? \writelines{2}
    \end{enumerate}
\end{enumerate}
\end{notesbox}

\begin{extensionbox}
\textbf{The golden ratio: the number that is its own reciprocal, plus one.} The
\textbf{golden ratio} is
\[
  \varphi = \frac{1+\sqrt{5}}{2}.
\]
It shows up wherever a shape can be divided so that the part relates to the whole the way the
smaller piece relates to the part. Everything interesting about $\varphi$ is invisible until you
rationalize.

\begin{center}
\begin{tikzpicture}[scale=2.4, line join=round]
  % golden rectangle, width phi = 1.618, height 1
  \draw[thick, forest, fill=forestbg] (0,0) rectangle (1.618,1);
  \draw[thick, goldacc] (1,0) -- (1,1);
  \node[charcoal] at (0.5,0.55) {\footnotesize square};
  \node[charcoal] at (0.5,0.38) {\footnotesize $1\times 1$};
  \node[charcoal, align=center] at (1.309,0.5) {\footnotesize what's\\[-2pt] \footnotesize left};
  % dimension labels
  \draw[<->, slate] (0,-0.13) -- (1.618,-0.13);
  \node[below, charcoal] at (0.809,-0.13) {\footnotesize $\varphi$};
  \draw[<->, slate] (-0.13,0) -- (-0.13,1);
  \node[left, charcoal] at (-0.13,0.5) {\footnotesize $1$};
  \draw[<->, slate] (1,1.13) -- (1.618,1.13);
  \node[above, charcoal] at (1.309,1.13) {\footnotesize $\varphi-1$};
\end{tikzpicture}
\end{center}

\begin{enumerate}[label=(\alph*), leftmargin=1.8em, itemsep=5pt]
  \item Using $\sqrt5\approx2.236068$, fill in the table to four decimal places.
        \begin{center}
        \renewcommand{\arraystretch}{1.4}
        \begin{tabular}{|c|c|c|}
        \hline
        $\varphi$ & $\dfrac{1}{\varphi}$ & $\varphi^{2}$ \\ \hline
        \blank{1.8cm} & \blank{1.8cm} & \blank{1.8cm} \\ \hline
        \end{tabular}
        \end{center}
        Something very strange is going on with those three decimals. Describe it in one sentence.
  \item \textbf{Prove the first strange thing:} $\dfrac{1}{\varphi}=\varphi-1$. Start from
        $\dfrac{1}{\varphi}=\dfrac{2}{1+\sqrt5}$, rationalize with the conjugate, and compare your
        result to $\varphi-1$.
  \item \textbf{Prove the second:} $\varphi^{2}=\varphi+1$. Square $\dfrac{1+\sqrt5}{2}$, simplify,
        and compare to $\varphi+1$ written over a common denominator.
  \item \textbf{The rectangle.} The figure above is $\varphi$ wide and $1$ tall. Cut off the
        $1\times1$ square. The leftover rectangle is $1$ tall and $\varphi-1$ wide, so its
        \emph{long side divided by its short side} is $\dfrac{1}{\varphi-1}$. Rationalize that
        expression and show it equals $\varphi$ --- the leftover rectangle has exactly the same
        proportions as the original.
  \item What does part (d) mean about repeating the cut on the leftover rectangle, and then again
        on \emph{its} leftover? Why is this one of the few facts about $\varphi$ you could not have
        seen without rationalizing?
\end{enumerate}
\writelines{6}
\end{extensionbox}

\begin{spiralbox}
\textbf{Coming next --- Lesson 5.4: Graphing Radical Functions \& Transformations.} The algebra of
radicals is now finished: you can simplify them, add them, multiply them, and clear them out of a
denominator. Next you go back to the \emph{pictures} from Lesson 5.0 and learn to move them around:
$y=a\sqrt{x-h}+k$ and $y=a\sqrt[3]{x-h}+k$, where the \textbf{endpoint} is the anchor that every
transformation moves, and $a$, $h$, and $k$ do exactly what they did to parabolas back in Unit 2.
You will read a graph and write its equation, and write an equation and predict its graph.
\end{spiralbox}

\end{document}
